% Part:sets-functions-relations
% Chapter: size-of-sets
% Section: zig-zag

\documentclass[../../../include/open-logic-section]{subfiles}

\begin{document}

\olfileid{sfr}{siz}{zigzag}

\olsection{Cantor 之字形方法}

\begin{explain}
我们已经考虑了一些“简单”的枚举。现在我们要考虑稍难一些的情况。考虑自然数对的集合\oliflabeldef{sfr:set:pai:sec}{，我们在\olref[set][pai]{sec}中如此定义：}{，其定义如下：}
\[
\Nat \times \Nat = \Setabs{\tuple{n,m}}{n,m \in \Nat}
\]
我们可以将这些有序对组织成一个\emph{阵列}，如下所示：
\[
\begin{array}{ c | c | c | c | c | c}
& \mathbf 0 & \mathbf 1 & \mathbf 2 & \mathbf 3 & \dots \\
\hline
\mathbf 0 & \tuple{0,0} & \tuple{0,1} & \tuple{0,2} & \tuple{0,3} & \dots \\
\hline
\mathbf 1 & \tuple{1,0} & \tuple{1,1} & \tuple{1,2} & \tuple{1,3} & \dots \\
\hline
\mathbf 2 & \tuple{2,0} & \tuple{2,1} & \tuple{2,2} & \tuple{2,3} & \dots \\
\hline
\mathbf 3 & \tuple{3,0} & \tuple{3,1} & \tuple{3,2} & \tuple{3,3} & \dots \\
\hline
\vdots & \vdots & \vdots & \vdots & \vdots & \ddots\\
\end{array}
\]
显然，$\Nat \times \Nat$ 中的每个有序对都会在这个阵列中恰好出现一次。确切地说，$\tuple{n,m}$ 将出现在第 $n$ 行第 $m$ 列。但是，我们如何将这样一个阵列的元素组织成一个“一维”列表呢？下面阵列中的模式展示了其中一种方法（当然，还有许多其他选择）：
\[
\begin{array}{ c | c | c | c | c | c | c}
& \mathbf 0 & \mathbf 1 & \mathbf 2 & \mathbf 3 & \mathbf 4 &\dots \\
\hline
\mathbf 0 & 0  & 1& 3 & 6& 10 &\ldots \\
\hline
\mathbf 1 &2 & 4& 7 & 11 & \dots &\ldots \\
\hline
\mathbf 2 & 5 & 8 & 12 & \ldots & \dots&\ldots \\
\hline
\mathbf 3 & 9 & 13 & \ldots & \ldots & \dots & \ldots \\
\hline
\mathbf 4 & 14 & \ldots & \ldots & \ldots & \dots & \ldots \\
\hline
\vdots & \vdots & \vdots & \vdots & \vdots&\ldots & \ddots\\
\end{array}
\]\noindent
这种模式称为\emph{Cantor 之字形方法}。它按如下方式枚举 $\Nat \times \Nat$：
\[
\tuple{0,0}, \tuple{0,1}, \tuple{1,0}, \tuple{0,2}, \tuple{1,1},
\tuple{2,0}, \tuple{0,3}, \tuple{1,2}, \tuple{2,1}, \tuple{3,0}, \dots
\]
由此我们得到如下结论：
\end{explain}

\begin{prop}\ollabel{natsquaredenumerable}
$\Nat \times \Nat$ 是可枚举的。
\end{prop}

\begin{proof}
令 $f \colon \Nat \to \Nat\times\Nat$ 将每个 $k \in \Nat$ 映射到有序对 $\tuple{n,m} \in \Nat \times \Nat$，其中 $k$ 是 Cantor 之字形阵列中第 $n$ 行第 $m$ 列的值。
\end{proof}

\begin{explain}
这个技巧可以很好地推广。例如，我们可以用它来枚举自然数有序三元组的集合，即：
\[
\Nat \times \Nat \times \Nat = \Setabs{\tuple{n,m,k}}{n,m,k \in \Nat}
\]
我们可以将 $\Nat \times \Nat \times \Nat$ 视为 $\Nat \times \Nat$ 与 $\Nat$ 的笛卡尔积，也就是说，
\[
\Nat^3 = (\Nat \times \Nat) \times \Nat =
\Setabs{\tuple{\tuple{n,m},k}}{n, m, k
  \in \Nat }
\]
因此，我们可以用一个阵列来枚举 $\Nat^3$，方法是在一个轴上标出 $\Nat$ 的枚举，在另一个轴上标出 $\Nat^2$ 的枚举：
\[
\begin{array}{ c | c | c | c | c | c}
& \mathbf 0 & \mathbf 1 & \mathbf 2 & \mathbf 3 & \dots \\
\hline
\mathbf{\tuple{0,0}} & \tuple{0,0,0} & \tuple{0,0,1} & \tuple{0,0,2} & \tuple{0,0,3} & \dots \\
\hline
\mathbf{\tuple{0,1}} & \tuple{0,1,0} & \tuple{0,1,1} & \tuple{0,1,2} & \tuple{0,1,3} & \dots \\
\hline
\mathbf{\tuple{1,0}} & \tuple{1,0,0} & \tuple{1,0,1} & \tuple{1,0,2} & \tuple{1,0,3} & \dots \\
\hline
\mathbf{\tuple{0,2}} & \tuple{0,2,0} & \tuple{0,2,1} & \tuple{0,2,2} & \tuple{0,2,3} & \dots\\
\hline
\vdots & \vdots & \vdots & \vdots & \vdots & \ddots \\
\end{array}
\]
这样，通过使用与 Cantor 之字形方法类似的方法，我们同样可以得到~$\Nat^3$ 的一个枚举。我们可以继续这个过程，对于任意自然数 $n$，得到 $\Nat^n$ 的枚举。所以，我们有：
\end{explain}

\begin{prop}
对于每个 $n \in \Nat$，$\Nat^n$ 是可枚举的。
\end{prop}

\begin{prob}\label{sfr:siz:zigzag:prob:posint-n}
证明对于每个 $n \in \Nat$，$(\PosInt)^n$ 是可枚举的。
\end{prob}

\begin{prob}\label{sfr:siz:zigzag:prob:posint-star}
  证明 $(\PosInt)^*$ 是可枚举的。你可以假设 \cref{sfr:siz:zigzag:prob:posint-n}。
\end{prob}

\end{document}